\documentclass[11pt,a4paper]{article}

\usepackage[utf8]{inputenc}
\usepackage[T1]{fontenc}
\usepackage{lmodern}
\usepackage[margin=1in]{geometry}
\usepackage{amsmath,amssymb}
\usepackage{enumitem}
\usepackage{booktabs}
\usepackage{fancyhdr}
\usepackage{xcolor}
\usepackage{mdframed}

% Gray solution boxes
\definecolor{solngray}{gray}{0.95}
\newmdenv[backgroundcolor=solngray,linewidth=0.5pt,linecolor=gray,innertopmargin=8pt,innerbottommargin=8pt]{solnbox}

\pagestyle{fancy}
\fancyhf{}
\fancyhead[L]{\textbf{Discrete Structures I}}
\fancyhead[R]{\textbf{Prelim Practice Exam B}}
\fancyfoot[C]{\thepage}
\renewcommand{\headrulewidth}{0.5pt}

\begin{document}

\begin{center}
{\LARGE\textbf{DISCRETE STRUCTURES I}}\\[5pt]
{\Large Prelim Practice Exam --- Version B}\\[5pt]
{\large\textit{With Complete Solutions}}\\[10pt]
\rule{0.8\textwidth}{0.5pt}\\[5pt]
\textbf{Time Allowed:} 1.5 hours \quad | \quad \textbf{Total Points:} 55
\end{center}

\vspace{10pt}

\noindent\textbf{Instructions:}
\begin{itemize}[noitemsep]
    \item Answer all questions. Show all work for full credit.
\end{itemize}

\vspace{15pt}

%==============================================================================
% PART I: MULTIPLE CHOICE
%==============================================================================
\section*{PART I: Multiple Choice (10 points)}
\textit{2 points each.}

\vspace{8pt}

\textbf{1.} The statement ``$p$ only if $q$'' is logically equivalent to:

\begin{enumerate}[label=(\Alph*),noitemsep]
    \item $q \to p$
    \item $p \to q$
    \item $p \land q$
    \item $p \leftrightarrow q$
\end{enumerate}

\begin{solnbox}
\textbf{Answer: (B)}

\textbf{Explanation:} ``$p$ only if $q$'' is one of the trickiest English phrases to translate!

\textbf{Key insight:} ``$p$ only if $q$'' means ``$p$ can be true ONLY when $q$ is true.'' In other words, if $p$ is true, then $q$ MUST be true.

This translates to: $p \to q$

\textbf{Common mistake:} Students often confuse this with $q \to p$. Remember:
\begin{itemize}[noitemsep]
    \item ``If $p$ then $q$'' = $p \to q$
    \item ``$p$ only if $q$'' = $p \to q$ (same!)
    \item ``$p$ if $q$'' = $q \to p$ (different!)
\end{itemize}
\end{solnbox}

\vspace{10pt}

\textbf{2.} Which of the following is a tautology?

\begin{enumerate}[label=(\Alph*),noitemsep]
    \item $p \land \neg p$
    \item $p \lor \neg p$
    \item $p \to q$
    \item $p \land q$
\end{enumerate}

\begin{solnbox}
\textbf{Answer: (B)}

\textbf{Explanation:}
\begin{itemize}[noitemsep]
    \item (A) $p \land \neg p$ is a \textbf{contradiction} --- always FALSE (a thing can't be both true and false)
    \item (B) $p \lor \neg p$ is a \textbf{tautology} --- always TRUE (Law of Excluded Middle)
    \item (C) $p \to q$ is a \textbf{contingency} --- depends on values of $p$ and $q$
    \item (D) $p \land q$ is a \textbf{contingency} --- depends on values
\end{itemize}

$p \lor \neg p$ says ``either $p$ is true or $p$ is false'' --- this is always true!
\end{solnbox}

\vspace{10pt}

\textbf{3.} If $P(x)$ denotes ``$x > 5$'' and the domain is $\{1, 3, 5, 7, 9\}$, what is the truth value of $\exists x\, P(x)$?

\begin{enumerate}[label=(\Alph*),noitemsep]
    \item True
    \item False
\end{enumerate}

\begin{solnbox}
\textbf{Answer: (A)}

\textbf{Explanation:} $\exists x\, P(x)$ asks: ``Does there exist at least one element in the domain where $x > 5$?''

Check each element:
\begin{itemize}[noitemsep]
    \item $x = 1$: $1 > 5$? No
    \item $x = 3$: $3 > 5$? No
    \item $x = 5$: $5 > 5$? No (not strictly greater)
    \item $x = 7$: $7 > 5$? \textbf{Yes!}
    \item $x = 9$: $9 > 5$? \textbf{Yes!}
\end{itemize}

Since we found at least one (actually two!) elements satisfying $P(x)$, the existential statement is \textbf{TRUE}.

\textbf{Remember:} For $\exists$, you only need ONE witness. For $\forall$, you need ALL to satisfy.
\end{solnbox}

\vspace{10pt}

\textbf{4.} What is the correct translation of ``All mathematicians are logical''?

\begin{enumerate}[label=(\Alph*),noitemsep]
    \item $\forall x\, (M(x) \land L(x))$
    \item $\forall x\, (M(x) \to L(x))$
    \item $\exists x\, (M(x) \to L(x))$
    \item $\exists x\, (M(x) \land L(x))$
\end{enumerate}

\begin{solnbox}
\textbf{Answer: (B)}

\textbf{Explanation:} ``All A are B'' ALWAYS uses implication, not conjunction!

$$\text{``All A are B''} \equiv \forall x\, (A(x) \to B(x))$$

\textbf{Why not conjunction?}
\begin{itemize}[noitemsep]
    \item $\forall x\, (M(x) \land L(x))$ would mean ``Everything is a mathematician AND logical'' --- this claims everyone/everything is a mathematician!
    \item $\forall x\, (M(x) \to L(x))$ means ``For anything, IF it's a mathematician, THEN it's logical'' --- only makes claims about mathematicians.
\end{itemize}

\textbf{Contrast with ``Some'':} ``Some A are B'' uses conjunction: $\exists x\, (A(x) \land B(x))$
\end{solnbox}

\vspace{10pt}

\textbf{5.} Which argument form is INVALID (a fallacy)?

\begin{enumerate}[label=(\Alph*),noitemsep]
    \item $p \to q$, $p$, therefore $q$
    \item $p \to q$, $\neg q$, therefore $\neg p$
    \item $p \to q$, $q$, therefore $p$
    \item $p \lor q$, $\neg p$, therefore $q$
\end{enumerate}

\begin{solnbox}
\textbf{Answer: (C)}

\textbf{Explanation:}
\begin{itemize}[noitemsep]
    \item (A) Modus Ponens --- VALID
    \item (B) Modus Tollens --- VALID
    \item (C) \textbf{Affirming the Consequent} --- INVALID (FALLACY!)
    \item (D) Disjunctive Syllogism --- VALID
\end{itemize}

\textbf{Why is (C) invalid?}

Consider: ``If it rains, the ground is wet. The ground is wet. Therefore, it rained.''

But wait! The ground could be wet because someone used a hose! Just because $q$ is true doesn't mean $p$ caused it.

\textbf{The two classic fallacies:}
\begin{enumerate}[noitemsep]
    \item Affirming the Consequent: $p \to q$, $q$ $\therefore p$ (WRONG!)
    \item Denying the Antecedent: $p \to q$, $\neg p$ $\therefore \neg q$ (WRONG!)
\end{enumerate}
\end{solnbox}

\newpage

%==============================================================================
% PART II: SHORT ANSWER
%==============================================================================
\section*{PART II: Short Answer (20 points)}

\textbf{Question 1.} (5 points) Construct a truth table for $(\neg p \lor q) \to (p \land q)$.

\begin{solnbox}
\textbf{Identify subexpressions:}
\begin{enumerate}[noitemsep]
    \item $\neg p$
    \item $\neg p \lor q$ (this is actually equivalent to $p \to q$!)
    \item $p \land q$
    \item Final: $(\neg p \lor q) \to (p \land q)$
\end{enumerate}

\begin{center}
\begin{tabular}{cc|c|cc|c}
\toprule
$p$ & $q$ & $\neg p$ & $\neg p \lor q$ & $p \land q$ & $(\neg p \lor q) \to (p \land q)$ \\
\midrule
T & T & F & T & T & T \\
T & F & F & F & F & T \\
F & T & T & T & F & F \\
F & F & T & T & F & F \\
\bottomrule
\end{tabular}
\end{center}

\textbf{Row-by-row calculation:}

\textbf{Row 1} ($p=T, q=T$):
\begin{itemize}[noitemsep]
    \item $\neg p = F$, so $\neg p \lor q = F \lor T = T$
    \item $p \land q = T \land T = T$
    \item $T \to T = T$ (true implies true = true)
\end{itemize}

\textbf{Row 3} ($p=F, q=T$):
\begin{itemize}[noitemsep]
    \item $\neg p = T$, so $\neg p \lor q = T \lor T = T$
    \item $p \land q = F \land T = F$
    \item $T \to F = F$ (true implies false = FALSE!)
\end{itemize}

\textbf{Observation:} This is a contingency (has both T and F).

\textbf{Interesting note:} $\neg p \lor q \equiv p \to q$, so we're really looking at $(p \to q) \to (p \land q)$.
\end{solnbox}

\vspace{15pt}

\textbf{Question 2.} (5 points) Translate into propositional logic:

\textit{``The system will shut down if there is a power failure and the backup generator is not working.''}

Let:
\begin{itemize}[noitemsep]
    \item $s$: ``The system will shut down''
    \item $p$: ``There is a power failure''
    \item $b$: ``The backup generator is working''
\end{itemize}

\begin{solnbox}
\textbf{Parsing the English:}

The structure is: ``[A] if [B]'' which means ``If [B] then [A]'' or $B \to A$.

So: ``System shuts down IF [condition]'' means [condition] $\to s$.

What's the condition? ``Power failure AND backup NOT working'' = $p \land \neg b$

\textbf{Answer:}
$$\boxed{(p \land \neg b) \to s}$$

\textbf{Alternative reading:} Some might interpret this as the condition for shutdown. Either way, the core translation is the same.

\textbf{Common mistake:} Don't confuse ``if'' and ``only if''!
\begin{itemize}[noitemsep]
    \item ``A if B'' = $B \to A$ (B is sufficient for A)
    \item ``A only if B'' = $A \to B$ (B is necessary for A)
\end{itemize}
\end{solnbox}

\vspace{15pt}

\textbf{Question 3.} (4 points) Determine the truth value. Domain: all positive integers.

\begin{enumerate}[label=(\roman*)]
    \item $\forall x \forall y\, (x \cdot y \geq 1)$
    \item $\exists x \exists y\, (x - y = 0)$
\end{enumerate}

\begin{solnbox}
\textbf{(i) $\forall x \forall y\, (x \cdot y \geq 1)$}

``For all positive integers $x$ and $y$, $x \cdot y \geq 1$.''

Since $x \geq 1$ and $y \geq 1$ for all positive integers, their product $x \cdot y \geq 1 \cdot 1 = 1$.

This is \textbf{TRUE}.

\rule{\linewidth}{0.3pt}

\textbf{(ii) $\exists x \exists y\, (x - y = 0)$}

``There exist positive integers $x$ and $y$ such that $x - y = 0$.''

$x - y = 0$ means $x = y$.

Can we find positive integers that are equal? Of course! $x = y = 1$ works. (Or $x = y = 5$, etc.)

This is \textbf{TRUE}.
\end{solnbox}

\vspace{15pt}

\textbf{Question 4.} (6 points) Let $C(x, y)$ mean ``$x$ has called $y$'' where the domain is all people with phones. Express:

\begin{enumerate}[label=(\roman*)]
    \item Someone has called everyone.
    \item Everyone has been called by someone.
    \item No one has called themselves.
\end{enumerate}

\begin{solnbox}
\textbf{(i) Someone has called everyone.}

``There exists a person $x$ such that for all people $y$, $x$ has called $y$.''

$$\boxed{\exists x \forall y\, C(x, y)}$$

\rule{\linewidth}{0.3pt}

\textbf{(ii) Everyone has been called by someone.}

``For every person $y$, there exists some person $x$ such that $x$ has called $y$.''

$$\boxed{\forall y \exists x\, C(x, y)}$$

\textbf{Key insight:} Compare (i) and (ii):
\begin{itemize}[noitemsep]
    \item (i) says ONE person made calls to everyone (same caller, all recipients)
    \item (ii) says everyone received a call from SOMEONE (different callers possible)
\end{itemize}

\rule{\linewidth}{0.3pt}

\textbf{(iii) No one has called themselves.}

``For all people $x$, $x$ has NOT called $x$.'' (No self-calls!)

$$\boxed{\forall x\, \neg C(x, x)}$$

Alternatively: $\neg \exists x\, C(x, x)$ (There does not exist someone who called themselves)
\end{solnbox}

\newpage

%==============================================================================
% PART III: PROBLEM SOLVING
%==============================================================================
\section*{PART III: Problem Solving (25 points)}

\textbf{Question 5.} (6 points) Identify the rule of inference:

\begin{enumerate}[label=(\roman*)]
    \item Sam is a computer science student. Therefore, Sam is a computer science student or Sam is an art major.
    \item If today is Tuesday, then I have a quiz. If I have a quiz, then I studied last night. Therefore, if today is Tuesday, then I studied last night.
    \item It is raining and it is cold. Therefore, it is raining.
\end{enumerate}

\begin{solnbox}
\textbf{(i)} Let $c$ = ``Sam is a CS student'', $a$ = ``Sam is an art major''

Structure: $c$, therefore $c \lor a$

This is \textbf{Addition}!

$$\frac{p}{\therefore p \lor q}$$

\textbf{Why it works:} If something is true, adding it to a disjunction with anything else is still true (since OR only needs one true part).

\rule{\linewidth}{0.3pt}

\textbf{(ii)} Let $t$ = ``today is Tuesday'', $q$ = ``I have a quiz'', $s$ = ``I studied last night''

Structure: $t \to q$, $q \to s$, therefore $t \to s$

This is \textbf{Hypothetical Syllogism}!

$$\frac{p \to q \quad q \to r}{\therefore p \to r}$$

\textbf{Why it works:} Chain of implications. Tuesday leads to quiz, quiz leads to studying, so Tuesday leads to studying.

\rule{\linewidth}{0.3pt}

\textbf{(iii)} Let $r$ = ``raining'', $c$ = ``cold''

Structure: $r \land c$, therefore $r$

This is \textbf{Simplification}!

$$\frac{p \land q}{\therefore p}$$

\textbf{Why it works:} If both parts are true (conjunction), then certainly each individual part is true.
\end{solnbox}

\vspace{15pt}

\textbf{Question 6.} (7 points) Prove: $(p \to q) \land (p \to r) \equiv p \to (q \land r)$

\begin{solnbox}
\textbf{Strategy:} Convert implications to disjunctions, then manipulate.

\begin{align*}
(p \to q) \land (p \to r) &\equiv (\neg p \lor q) \land (\neg p \lor r) && \text{[Implication Law, twice]} \\
&\equiv \neg p \lor (q \land r) && \text{[Distributive Law (reverse)]} \\
&\equiv p \to (q \land r) && \text{[Implication Law]}
\end{align*}

$$\boxed{(p \to q) \land (p \to r) \equiv p \to (q \land r)}$$

\textbf{Distributive Law clarification:}

The Distributive Law says: $A \lor (B \land C) \equiv (A \lor B) \land (A \lor C)$

We used it in reverse: $(\neg p \lor q) \land (\neg p \lor r) \equiv \neg p \lor (q \land r)$

Here $A = \neg p$, $B = q$, $C = r$.

\textbf{Intuitive meaning:} If $p$ implies $q$ AND $p$ implies $r$, that's the same as saying $p$ implies ``$q$ and $r$''. If $p$ forces two separate conclusions, then it forces both together.
\end{solnbox}

\vspace{15pt}

\textbf{Question 7.} (4 points) Negate and simplify:

$$\exists x \forall y \exists z\, (P(x, y) \to Q(y, z))$$

\begin{solnbox}
\textbf{Step 1:} Push negation through quantifiers, flipping each one.

\begin{align*}
&\neg(\exists x \forall y \exists z\, (P(x,y) \to Q(y,z))) \\
&\equiv \forall x\, \neg(\forall y \exists z\, (P(x,y) \to Q(y,z))) && \text{[$\neg\exists \to \forall\neg$]} \\
&\equiv \forall x \exists y\, \neg(\exists z\, (P(x,y) \to Q(y,z))) && \text{[$\neg\forall \to \exists\neg$]} \\
&\equiv \forall x \exists y \forall z\, \neg(P(x,y) \to Q(y,z)) && \text{[$\neg\exists \to \forall\neg$]}
\end{align*}

\textbf{Step 2:} Simplify the negated implication.

Recall: $\neg(A \to B) \equiv A \land \neg B$

So: $\neg(P(x,y) \to Q(y,z)) \equiv P(x,y) \land \neg Q(y,z)$

\textbf{Final answer:}
$$\boxed{\forall x \exists y \forall z\, (P(x,y) \land \neg Q(y,z))}$$
\end{solnbox}

\vspace{15pt}

\textbf{Question 8.} (8 points) Prove valid:

\textbf{Premises:}
\begin{enumerate}[noitemsep]
    \item $\forall x\, (P(x) \to Q(x))$
    \item $\forall x\, (Q(x) \to R(x))$
    \item $P(a)$
\end{enumerate}

\textbf{Conclusion:} $R(a)$

\begin{solnbox}
\textbf{Strategy:} Use Universal Instantiation to get specific statements about $a$, then chain them.

\begin{center}
\begin{tabular}{c|l|l}
\toprule
\textbf{Step} & \textbf{Statement} & \textbf{Reason} \\
\midrule
1 & $\forall x\, (P(x) \to Q(x))$ & Premise \\
2 & $P(a) \to Q(a)$ & Universal Instantiation (1) \\
3 & $\forall x\, (Q(x) \to R(x))$ & Premise \\
4 & $Q(a) \to R(a)$ & Universal Instantiation (3) \\
5 & $P(a) \to R(a)$ & Hypothetical Syllogism (2, 4) \\
6 & $P(a)$ & Premise \\
7 & $R(a)$ & Modus Ponens (5, 6) \\
\bottomrule
\end{tabular}
\end{center}

$$\boxed{\therefore R(a)}$$

\textbf{Explanation:}
\begin{itemize}[noitemsep]
    \item Steps 1-2: The universal statement ``for all $x$, $P(x) \to Q(x)$'' applies to any specific element. We pick $a$.
    \item Steps 3-4: Similarly instantiate the second universal statement for $a$.
    \item Step 5: Chain the two implications: $P(a) \to Q(a)$ and $Q(a) \to R(a)$ gives $P(a) \to R(a)$.
    \item Steps 6-7: We HAVE $P(a)$ (premise 3), and we have $P(a) \to R(a)$, so by Modus Ponens, $R(a)$.
\end{itemize}

This is a classic ``chain of implications with universal statements'' problem!
\end{solnbox}

\vspace{20pt}

\begin{center}
\rule{0.5\textwidth}{0.5pt}\\[5pt]
\textbf{END OF EXAM}\\[3pt]
\textit{Total: 55 points}
\end{center}

\end{document}
